\renewcommand{\Anno}{2000}

\renewcommand{\NumeroProblema}{2}

\renewcommand{\Testo}{Nel centro di una piazza sono state sospese due
sfere.  La prima ha raggio $r$ ed il suo centro sta ad altezza $h$, la seconda ha raggio $R$ ed il suo centro sta ad altezza $h+d$ (si supponga $h>r$, $R>r$, $d>R+r$).  La retta congiungente i due centri è perpendicolare alla piazza.  Durante la notte, la seconda sfera si illumina, diffondendo la sua luce in tutte le direzioni.
	
\begin{enumerate}
	\item Determinare, in funzione di $r$, $R$, $h$, $d$, l'area della parte di piazza che rimane in ombra (costituita dai punti
	guardando dai quali la seconda sfera è \emph{totalmente}
	nascosta dietro la prima).
	
	\item Determinare, in funzione di $r$, $R$, $h$, $d$, l'area della parte di piazza che rimane in penombra (costituita dai punti
	guardando dai quali la seconda sfera è \emph{parzialmente}
	nascosta dietro la prima).
	
	\item Discutere il caso limite $d\to +\infty$ (cioè dire cosa succede alla zona d'ombra e di penombra quando $d$ è molto grande rispetto ad $r$, $R$, $h$).
\end{enumerate} }

\renewcommand{\Risposta}{\begin{enumerate} 
	\item[(a)] Se $rd\geq h(R-r)$, la zona d'ombra è un cerchio di area
	$$\frac{(rd-h(R+r))^{2}}{d^{2}-(R+r)^{2}}\cdot \pi;$$
	se $rd<h(R-r)$, la zona d'ombra è vuota.
	
	\item[(b)] Se $rd\geq h(R-r)$, la zona di penombra è una corona circolare di area
	$$\frac{4rdh(R+r)}{d^{2}-(R+r)^{2}}\cdot \pi;$$
	se $rd<h(R-r)$, la zona di penombra è un cerchio di area
	$$\frac{(rd+h(R+r))^{2}}{d^{2}-(R+r)^{2}}\cdot \pi.$$
	
	\item[(c)] Quando $d\to+\infty$ il raggio della zona d'ombra ed il raggio esterno della zona di penombra tendono ad $r$.  Pertanto la zona d'ombra cresce fino a diventare un cerchio di raggio $r$ e la zona di penombra è una corona circolare sempre più stretta fino a sparire al limite.
\end{enumerate}}

\renewcommand{\Aiutino}{Disegnare la sezione della figura mediante il
piano perpendicolare a terra e passante per i centri delle due sfere. Tracciare quindi le 4 rette tangenti ad entrambi i cerchi e determinare, mediante opportune similitudini, le loro intersezioni con
la retta di terra.}

\renewcommand{\SoluzioneUfficiale}{
}

\renewcommand{\SoluzioniAlternative}{
}

\renewcommand{\Osservazioni}{
}

\renewcommand{\Commenti}{
}
% Aggiunta redazionale: la soluzione preesistente non viene modificata.
\providecommand{\SoluzioneRedazionale}{}
\renewcommand{\SoluzioneRedazionale}{
% SSSUP-BOZZA-AI: SoluzioneRedazionale; autore=ChatGPT; data=2026-11-08; stato=DA-VERIFICARE
% SSSUP-PROVENIENZA: AI: elaborazione nuova dal testo
Poniamo $H=h+d$. Per simmetria, ombra e penombra sono concentriche rispetto al punto della piazza sotto i due centri. Introduciamo
\[
\rho_u=\max\left\{0,\frac{rd-h(R-r)}{\sqrt{d^2-(R-r)^2}}\right\},\qquad
\rho_p=\frac{rd+h(R+r)}{\sqrt{d^2-(R+r)^2}}.
\]
I denominatori sono positivi grazie a $d>R+r$.

\paragraph{Geometria delle tangenti.}
Consideriamo la sezione verticale per i centri. Le tangenti comuni esterne ai due cerchi delimitano il cono d'ombra. Esse si incontrano sull'asse a quota
\[
z_u=h-\frac{rd}{R-r},
\]
e formano con l'asse un angolo $\gamma$ tale che $\sin\gamma=(R-r)/d$. Se $z_u\leq0$, l'intersezione del cono con il terreno ha raggio $(-z_u)\tan\gamma$, che è precisamente $\rho_u$. Se $z_u>0$, il cono termina sopra il terreno e non vi è una regione di ombra totale.

Le tangenti comuni interne delimitano la regione nella quale almeno parte della sfera luminosa è occultata. Il loro punto d'incontro ha quota
\[
z_p=h+\frac{rd}{R+r},
\]
e l'angolo con l'asse soddisfa $\sin\eta=(R+r)/d$. Il raggio esterno della zona interessata è $z_p\tan\eta=\rho_p$.

\paragraph{Punti (a) e (b).}
Trascurando i bordi, che hanno area nulla, le aree richieste sono quindi
\[
\operatorname{Area}(\text{ombra})=\pi\rho_u^2,
\qquad
\operatorname{Area}(\text{penombra})=\pi(\rho_p^2-\rho_u^2).
\]
Se $rd<h(R-r)$, la prima regione è vuota e la penombra è il disco di raggio $\rho_p$. Quando $rd=h(R-r)$, l'ombra si riduce al punto sull'asse e ha area zero.

\paragraph{Punto (c).}
Per $d\to+\infty$ entrambi i raggi tendono a $r$. L'area d'ombra tende dunque a $\pi r^2$, mentre l'area di penombra tende a zero.
}
